\section{Methodology} \label{sec:method}

\subsection{Overview}

In this work, we propose Event-B Agent, an LLM-powered end-to-end framework for formal model synthesis and repair.
\autoref{fig:Overview} shows the overall workflow of Event-B Agent, the model \& proof repair component is elaborated in \autoref{fig:modelrepair_v2}, and \autoref{fig:running_example} illustrates the minimum-searching model generated by Event-B Agent.
System description in natural language is provided to Event-B Agent to generate a refinement strategy, which distributes requirements to various refinement steps.
In this process, gluing invariants between two refinement steps are summarized in natural language.
During model synthesis, the distributed requirements and gluing invariants are formalized into a model.
The correctness of the model is then verified through model checking and theorem proving, and repairs are applied accordingly.
This procedure is repeated for each refinement step until all proof obligations are discharged or the trial limit is reached.

Event-B Agent adopts a neurosymbolic design, where semantic tasks such as refinement planning, model synthesis, and repair are delegated to specialized LLMs, while deterministic components handle the rest.
Model checkers, SMT solvers, and theorem provers verify synthesized models, pattern matching identifies candidate repair rules from the proof state, and atomic repair functions apply updates to models and proofs.
This separation combines the flexibility of LLMs with the reliability of symbolic reasoning, achieving both semantic versatility and soundness.

\begin{figure}[t]
\includegraphics[width=0.9\linewidth]{figures/overview_v2.drawio.pdf}
\caption{Overview of Event-B Agent. From natural-language requirements ($REQ$), the system plans a refinement strategy to guide incremental model synthesis, with verification and repair at each refinement level.}
\label{fig:Overview}
\centering
\end{figure}

\subsection{Refinement Strategy Planning}

Automated reasoning over realistic software systems is challenging due to the rapid growth of the search space as model complexity increases.
Refinement offers a complementary approach by establishing properties in simpler abstract models, where proofs are easier to discharge due to abstraction, and preserving them in more detailed concrete models through refinement proof obligations.
This ensures that proof effort scales with model complexity while preserving correctness.

On this basis, Event-B Agent utilizes LLMs to automatically plan refinement strategies from natural language requirements. 
Refinement introduces additional detail in a disciplined manner, guided by gluing invariants that relate abstract and concrete models to maintain consistency.

\subsubsection{Refinement} \label{sec:refinement}

Having introduced the role of refinement, we now formalize this notion.
Let $M_A$ denote an abstract model and $M_C$ a concrete refinement of $M_A$.
Let $REQ_{M}$ denote the set of requirements encoded in model $M$.
Refinement preserves requirements of the abstract model $REQ_{M_A}$ while introducing additional ones $REQ_{\mathit{diff}}$, i.e.,
\(
REQ_{M_A} \subset REQ_{M_C}
\wedge
REQ_{\mathit{diff}} \;\triangleq\; REQ_{M_C} \setminus REQ_{M_A}
\).

%% Editted version
% Next, we describe the formal definitions for proof obligations specific to refinement \delete{with the definitions introduced in Section~\ref{sec:preliminary}}.
% \edit{In addition to the definitions in \autoref{fig:grammar}, we use $H_{w,y}$ and $W_{x,w,y}$ to represent the guards and witnesses of a concrete event $e_C$ in $M_C$, where $H_{w,y} ::= \{pred_{s \cup c \cup w \cup y}\}$ and $W_{x,w,y} ::= \{pred_{s \cup c \cup x \cup w \cup y}\}$, and $w$, $y$ represent the set of variables in $M_C$ and set of parameters of $e_C$ disjoint from $v$ and $x$ respectively.}

% \delete{Next, we describe the formal definitions for proof obligations specific to refinement with the definitions introduced in Section~\ref{sec:preliminary}.}

% \noindent \edit{(Removed the POs.)}

% \begin{itemize}
%     \item $\mathit{GRD}$ (Guard Strengthening):
%     {\footnotesize \begin{equation}
%     \forall \; g_{v,x} \in G_{v,x} \; \cdot \; \Gamma, H_{w,y}, W_{x,w,y} \;\vdash\; g_{v,x}.
%     \end{equation}}

%     \item $\mathit{MRG}$ (Guard Merging):
%     {\footnotesize \begin{equation}
%     \Gamma, H_{w,y} \;\vdash\; G1_{v,x} \lor G2_{v,x}.
%     \end{equation}}

%     \item $\mathit{SIM}$ (Simulation):
%     {\footnotesize \begin{equation}
%     \Gamma, H_{w,y}, W_{x,w,y,w'}, BA_{C\_w,y,w'} \;\vdash\; BA_{A\_v,x,v'}.
%     \end{equation}}
% \end{itemize}
% \edit{(Updated the POs.)}

% \delete{When an abstract event is refined by a concrete one, the guard strengthening (GRD) PO ensures that the concrete guards are stronger than their abstract counterparts, preserving properties proven at the abstract level.
% If two abstract events need to be merged into a single concrete event, the guard merging (MRG) PO requires the concrete guards to be stronger than the disjunction of the abstract guards.
% Finally, when refinement introduces new variables or event parameters, the simulation (SIM) PO guarantees that, given the concrete guards and witnesses relating $x$ and $y$, every abstract action is simulated by the concrete actions. 
% For example, if $y=\varnothing \land w=\varnothing$, each abstract action appears in $BA_C(\varnothing,\varnothing)$.}

The preservation of proven properties is guaranteed through gluing invariants and the corresponding refinement POs.
For instance, the refined events are required to have stronger guards than the abstract events, and their actions must simulate the abstract counterparts.
With these guarding POs, Event-B Agent ensures refinement correctness, as illustrated in \autoref{fig:running_example}.

\subsubsection{Gluing Invariant} \label{sec:gluinginv}

A gluing invariant is a predicate over both abstract variables and concrete variables, specifying how states in the concrete model correspond to states in the abstract model.
\autoref{fig:running_example} includes a gluing invariant between $M_{abstract}$ and $M_{concrete}$ generated by Event-B Agent.
$M_{concrete}$ introduces a variable $i$ to represent the current index during search, and a boolean variable $searching$ to represent whether the searching procedure is ongoing.
The gluing invariant relates the abstract variable $j$ to the two concrete variables $i$ and $searching$, and specifies that during searching, $f(j)$ is always the minimum value among the explored domain $0..(i-1)$.
The model synthesis step later formalizes it into $searching=TRUE \land 1\leq i \land i \leq n \Rightarrow j \in 0..(i-1) \land f(j) = min(f[0..(i-1)])$.

Gluing invariants are added to the model $M$ and play a critical role in discharging refinement POs.
In Event-B Agent, candidate gluing invariants are first proposed by the refinement strategy planning LLM and then formalized during model synthesis.
Since LLM-based generation offers no correctness guarantee, gluing invariants are validated in two steps: (1) counterexample checking and contradiction detection with the model checker, and (2) attempting proofs relevant to the gluing invariants before other proofs.
Only after no counterexamples are found and the proofs succeed, the gluing invariants will be accepted into $M$ for subsequent reasoning.

\subsubsection{Refinement Strategy}

Refinement proceeds inductively, distributes requirements across successive steps, and preserves established properties.
The first stage of Event-B Agent is to plan a refinement strategy that decides how requirements are allocated across steps and how successive models are linked by gluing invariants.
Once the strategy is set, the model synthesis and repair components construct the models, verify the requirements, and apply corrections if needed.

Concretely, the refinement strategy planning step of Event-B Agent partitions $REQ$ into $n$ disjoint subsets $REQ_{M_i}$ such that $REQ=\bigcup_{i=1}^{n} REQ_{M_i}$.
$REQ_{M_1}$ includes the requirements for the initial abstract model, $REQ_{M_2}$ contains the additional requirements for the second model, etc.
Between steps $i-1$ and $i$, the refinement strategy planning LLM also proposes a set of gluing invariants $\mathcal{I}_{g_i}$, which are formalized and added into $M_i$ during model synthesis.
Thus, the additional requirements for step $i$ are $REQ_{M_i} \cup \mathcal{I}_{g_i}$, where $REQ_{M_i}$ specifies system requirements and $\mathcal{I}_{g_i}$ aims to preserve all previously satisfied requirements $\bigcup_{j=1}^{i-1} REQ_{M_j} \cup \mathcal{I}_{g_j}$, where $\mathcal{I}_{g_1}=\varnothing$.

In the minimum-searching example, Event-B Agent creates a refinement strategy of 2 steps, $REQ_{M_1} = \{FUN\text{--1}\}$ and $REQ_{M_2} = \{FUN\text{--2}, FUN\text{--3}, FUN\text{--4}, FUN\text{--5}, FUN\text{--6}\}$.
Additionally, as shown in \autoref{fig:running_example}, a set of gluing invariants $\mathcal{I}_{g_2} = \{inv_{g_2}\}$ is generated to capture the relation between refinement steps 1 and 2.
At this stage, the requirements in $REQ_{M_1}$, $REQ_{M_2}$, and $\mathcal{I}_{g_2}$ are in natural language form and will be formalized during the model synthesis step.

Formally, the correctness of a refined model $M_i$ with respect to all requirements up to step $i$ is defined inductively as:
% {\footnotesize \begin{equation}
% \begin{split}
%     Correct(M_i) \;\triangleq\; 
%     \Bigl( \bigwedge_{j=1}^{i-1} Correct(M_j) \Bigr) 
%     \land \Bigl( REQ_i \cup \mathcal{I}_{g_i} \subseteq \mathcal{I}(M_i) \cup \mathcal{V}(M_i) \cup \mathcal{A}(M_i) \cup \mathcal{E}(M_i) \Bigr) \\
%     \land \Bigl(\forall po \in PO_{M_i} \cdot Discharged(M_i, po) \Bigr)
% \end{split}
% \label{eq:refinement_correctness}
% \end{equation}}
% \edit{(Added \(\cup \mathcal{E}(M_i)\) in ~\autoref{eq:refinement_correctness}).}
{\footnotesize \begin{equation}
    Correct(M_i) \;\triangleq\; 
    \Bigl(\bigwedge_{j=1}^{i-1} Correct(M_j)\Bigr) 
    \land \Bigl(\bigwedge_{req \; \in \;REQ_{M_i} \cup \mathcal{I}_{g_i}} Covered(M_i, req)\Bigr) 
    \land \Bigl(\bigwedge_{po \in PO_{M_i}} Discharged(M_i, po)\Bigr)
\label{eq:refinement_correctness}
\end{equation}}
% \edit{(Simplified this equation).}

This planning stage provides the blueprint for synthesis and repair: each step has well-defined requirements and invariants, while correctness is ensured and preserved through proof obligations.


\subsection{Model Synthesis} \label{sec:modelsynthesis}

For each refinement step $i$ with requirements $REQ_{M_i} \cup \mathcal{I}_{g_i}$, Event-B Agent synthesizes a formal model that captures the specified requirements.
The model must be syntactically correct and free of compilation errors so that formal verification tools can be applied effectively.

\noindent
\textbf{Schema-Guided Formalization.} 
Simple prompt-based generation frequently produces ill-formed code, including undeclared variables, malformed invariants, and syntactic errors that are not parsable, particularly in the synthesis of formal specifications where training data is sparse.
To address this, we design a JSON schema that encodes the grammar of Event-B shown in \autoref{fig:grammar}.
The schema encodes structural constraints, for example, every machine must contain at least one event, and may include variables, invariants, and variants.
Because the schema is language-agnostic, it can be readily adapted to other formal specification languages, making our approach generalizable.

\noindent
\textbf{Synthesis and Repair Loop for Well-Formedness.}
The model synthesis LLM first generates a candidate model according to the JSON schema.
While the schema eliminates most syntax errors, it cannot enforce typing constraints, so issues such as undeclared variables, type mismatches, and invariants referencing undefined sets may still arise.
The candidate model is parsed and compiled into Event-B code.
Compilation errors are fed back to the LLM, together with the current model, for iterative repair to ensure well-formedness.
This synthesis–repair loop ensures that models passed to verification tools are syntactically valid and type-correct, providing a reliable foundation for subsequent refinement and semantic repair.

\noindent \textbf{Ensuring Refinement Soundness.}
While synthesizing a refined model, the requirements at the current refinement level and the previous abstract model are provided to the LLM as context.
The refinement relation is preserved through refinement POs after generating the refined model.
In our experiments, the refinement PO discharge rate is collected to evaluate the refinement correctness.

\subsection{Model \& Proof Repair}
\label{sec:repair}
A key contribution of our work is the model \& proof repair component, which iteratively verifies models, diagnoses inaccuracies, and updates both models and proofs.
An overview of this component is shown in \autoref{fig:modelrepair_v2}.

\begin{figure}[t]
\includegraphics[width=\linewidth]{figures/modelfixing_v2.drawio.pdf}
\caption{Workflow of the model \& proof repair component. The model is \textcircled{1} checked by a model checker (bounded) and repaired, then \textcircled{2} repaired guided by unbounded theorem proving.
}
\label{fig:modelrepair_v2}
\centering
\end{figure}

\subsubsection{Formal Verification and Repair}

Whenever a new model is synthesized, it is first verified by a model checker to detect violations of invariants, including both requirement and gluing invariants (Step \textcircled{1} in \autoref{fig:modelrepair_v2}).
The model repair LLM suggests suitable bounds, within which the model checker searches for states that violate invariants $\mathcal{I}$.
If a counterexample trace is found, it indicates an invariant violation. This trace is then passed to the model repair LLM as feedback for correcting the model to resolve the violation within the given bounds.

While model checking is effective for detecting invariant violations within the given bounds, it cannot guarantee full correctness.
This is because bounded analysis may miss unreachable states or behaviors outside the explored bounds.
To cover all executions, theorem proving is applied in addition.
After proof obligations are generated, automated provers and constraint solvers attempt to discharge them under the relevant hypotheses extracted from the model.
When proofs fail, the cause may range from modeling inaccuracies or limitations of the automated reasoning process.
To address this, the model \& proof repair component analyzes the proof state and applies targeted fix strategies to either the model or its proofs (Step \textcircled{2} in \autoref{fig:modelrepair_v2}).

\begin{table}[h!]

\footnotesize
\centering

% \setlength{\tabcolsep}{2pt}
\begin{tabular}{p{0\linewidth} p{0.13\linewidth} p{0.23\linewidth} p{0.53\linewidth}}
\toprule
\textbf{} & \textbf{Category} & \textbf{Condition} & \textbf{Rules} \\ 

\midrule
1 & \begin{tabular}[c]{@{}l@{}}
Contradictory \\
Goal
\end{tabular} & 
\begin{tabular}[c]{@{}l@{}}
Goal is $\bot$
\end{tabular} & 
\begin{tabular}[c]{@{}l@{}}
1. PO is about an invariant $\rightarrow$ update the invariant to eliminate \\
contradictions \\
2. PO is about an event $\rightarrow$ add or update actions or guards to \\
eliminate contradictions
\end{tabular} \\

\midrule
2 & \begin{tabular}[c]{@{}l@{}}
True by \\
Definition 
\end{tabular} & 
\begin{tabular}[c]{@{}l@{}}
Goal matches definitional \\
pattern, 
e.g. $x=y$, $x \in S$ etc.
\end{tabular} & 
\begin{tabular}[c]{@{}l@{}}
1. Goal holds by definition $\rightarrow$ apply the corresponding tactics \\
2. (a) Rule 1 doesn't update $\mathcal{P}$, or (b) Goal is false $\rightarrow$ Similar rules\\
as \textbf{Category 1}
\end{tabular} \\

\midrule
3 & Existential Goal & 
\begin{tabular}[c]{@{}l@{}}
Goal is an $\exists$ predicate, \\
e.g. $\exists x \cdot <predicate>$
\end{tabular} & 
\begin{tabular}[c]{@{}l@{}}
1. $x$ is a variable $\rightarrow$ add or update actions to define $x$ \\
2. $x$ is a constant $\rightarrow$ add hypothesis as an axiom in the context to \\
establish its existence
\end{tabular} \\

\midrule
4 & 
\begin{tabular}[c]{@{}l@{}}
Equality PO
\end{tabular} & 
\begin{tabular}[c]{@{}l@{}}
PO type is Equality, e.g. goal \\
$G$ is $x=y$, where $x$, $y$ are \\
from machines at different \\
refinement steps
\end{tabular} & 
\begin{tabular}[c]{@{}l@{}}
1. If there is an indirect relation $H$ (e.g. $f(x)=f(y)$), add axiom \\
of the form $H \Rightarrow G$ if suitable.
\end{tabular} \\

\midrule
5 & \begin{tabular}[c]{@{}l@{}}
Well-\\
Definedness 
\end{tabular} & 
\begin{tabular}[c]{@{}l@{}}
PO type is Well-Definedness \\
e.g. ensure $y \neq 0$ for $x/y$, \\
ensure $x \in dom(f)$ for $f(x)$
\end{tabular} & 
\begin{tabular}[c]{@{}l@{}}
1. PO is about guard $\rightarrow$ strengthen the guard based on the \\
definition of well-definedness \\
2. PO is about invariant $\rightarrow$ strengthen the invariant based on the \\
definition of well-definedness
\end{tabular} \\

\midrule
6 & 
\begin{tabular}[c]{@{}l@{}}
Quantified \\
Invariant \\
Preservation
\end{tabular} & 
\begin{tabular}[c]{@{}l@{}}
PO is Invariant Preservation \\
about a quantified invariant
\end{tabular} & 
\begin{tabular}[c]{@{}l@{}}
1. Universal invariant $\rightarrow$ Find quantified hypothesis \\
$\forall x \cdot H(x) \Rightarrow G(x)$, where $G(x)$ is similar to proof goal $G(x')$. \\
Instantiate the hypothesis with $x'$. \\
2. Existential invariant $\exists x \cdot P(x)$ $\rightarrow$ Similar 
rules as \textbf{Category 3}
\end{tabular} \\

\midrule
7 & 
\begin{tabular}[c]{@{}l@{}}
Uninstantiated \\
Hypothesis
\end{tabular} & 
\begin{tabular}[c]{@{}l@{}}
Previous repairs add \\
quantified hypotheses that \\
are not yet instantiated
\end{tabular} & 
\begin{tabular}[c]{@{}l@{}}
1. Quantified hypothesis added but not 
instantiated $\rightarrow$ suggest \\
suitable instantiation values
\end{tabular} \\

\bottomrule
\end{tabular}

\caption{Categories of proof states and the corresponding repair rules supported by Event-B Agent.}
\label{tab:rule-categories}

\end{table}


\subsubsection{Repair Rules Recommendation} \label{sec:proofrules}
To improve the reasoning capability of the fix strategy decision LLM, we design a novel repair rules recommendation component that suggests suitable rules for repairing the given model $M$ and a proof $\pi$ about the model, and provides them as guidance to the LLM.
We observe that characteristic patterns emerge across model $M$, proof $\pi$, the proof obligation type, and their corresponding historical repairs $R(M, \pi)$.

To illustrate how repair rules are recommended, consider the case of universally quantified hypotheses. 
In interactive theorem proving, such hypotheses become useful only after instantiation with concrete terms.
Although SMT solvers perform heuristic instantiation, this process is incomplete and often fails in Event-B proofs.
Thus, after adding a universal predicate as a lemma in $\pi$, the next step is typically to instantiate it. 
If the repair history $R(M, \pi)$ indicates that a universal lemma has been added but no instantiations appear, and the child nodes in the proof tree remain open, the following rule is recommended:
\textit{``If a universal lemma has been added to a node in the proof tree, but its child nodes remain unclosed, suggest instantiation values for the lemma''}.  
In this example, the first action modifies $M$ by adding the universal lemma as a reusable axiom in the context, while the second action instantiates it in $\pi$ if SMT solvers do not already discharge the proof.

Beyond this example, our system supports a total of seven categories of rules, summarized in \autoref{tab:rule-categories}. 
The ``Condition'' column specifies how each category is determined from the current proof state, including information from the model $M$, proof $\pi$, repair history $R(M, \pi)$, and proof obligation type.
Within each category, multiple concrete repairs may apply, and these sometimes cannot be determined by pattern matching. 
Instead, the fix strategy decision LLM analyzes the scenario, considers the recommended rules, and attempts repairs. 
The ``Rules'' column lists a non-exhaustive set of simplified rules under various scenarios.
These categories and rules were derived empirically based on our experience in manually discharging proofs.
While not exhaustive, they capture recurring patterns observed across a wide range of models.
If the current proof state does not fall into one of the known categories, the LLM falls back on a set of general default rules.

\autoref{fig:modelrepair_v2} shows an example of model-proof repair with the default rules.
$FUN$--3 in \autoref{tab:example} specifies a property about initial values of variables, formalized as $inv: \; n>0 \land searching=TRUE \land i=1 \Rightarrow j=0$.
The PO $find\_not\_better/FUN\_3$ fails because an event instance with pre-state $i=0$ and post-state $i'=1$ is admitted, yielding a post-state invariant $inv'$ that requires $j=0$, which cannot be derived from the proof context.
However, such a state is unreachable in the actual execution of the algorithm, where $i \geq 1$ is maintained (cf. \autoref{fig:running_example}).
As the proof state does not match any rule category in \autoref{tab:rule-categories}, the repair rules recommendation component falls back to default rules, whose simplified forms are shown as ``Suitable Repair Rules'' in \autoref{fig:modelrepair_v2}.
These rules guide the later steps to focus on the inductive relationship between $inv$ and $inv'$ (Rule 1), and add relevant hypotheses to $inv'$ when needed (Rule 2).

\subsubsection{Fix Strategy Selection and Execution}\label{sec:fix_selection}

The repair rules introduced in Section~\ref{sec:proofrules} are used to guide the fix strategy decision LLM, which may hallucinate during repair.
To mitigate this issue, Event-B Agent restricts modifications to a library of atomic repair functions (e.g., strengthen an invariant) that correspond to a set of fix strategies to update models and proofs.
The functions are described as ``atomic'' because they represent the smallest modification units.
After the repair rules are obtained, the LLM selects an atomic function and proposes the corresponding parameters, and the proof state is updated after executing the function.
This iterative process continues until the proof is discharged or the trial limit is reached.
The commonly used fix strategies are summarized in \autoref{fig:distribution}, which are grouped into 4 categories:
\begin{itemize}
    \item \textbf{Model modification strategies}: This group of strategies corresponds to atomic functions such as adding or updating invariants, strengthening guards, and modifying actions;
    \item \textbf{Proof modification strategies}: These strategies correspond to atomic functions such as instantiating quantified hypotheses and unfolding definitional equalities;
    \item \textbf{Joint model–proof modification strategies}: Operations that simultaneously modify the model and the selected proof, e.g., introducing a hypothesis as a context axiom while injecting it into the proof context, thereby repairing both in tandem;
    \item \textbf{Information Retrieval strategies}: Operations such as invoking the model checker to obtain guidance for subsequent repair steps, without directly modifying the model or proof.
\end{itemize}

To summarize, the fix strategy decision LLM performs two complementary tasks: 
(1) selecting the appropriate function based on the current proof state and recommended repair rules; and 
(2) proposing the values of function parameters, which specify the repair contents.
For instance, in \autoref{fig:modelrepair_v2}, the selected function is $strengthenGuard$, with parameters $guard\_label$ and $guard\_content$.
The repair adds $grd7:1 \leq i$, inferred from the inductive relationship between $inv$ and $inv'$.


\subsubsection{Soundness of Repair.}
In our method, repairs aimed at discharging POs are generated jointly by the repair rule recommendation component and the fix strategy decision LLM.
Based on pattern matching over the proof state, the former retrieves repair rules to construct prompts for the latter.
A sequence of repairs is accepted only if their combined effect discharges the target PO.
Once discharged, SMT solvers and theorem provers guarantee its validity.
After each model modification, all proofs are replayed to ensure that any invalidated proofs are not considered successful.
As a result, soundness is ensured by the verification pipeline rather than the repair rules.
The rules provide heuristic guidance by suggesting common proof actions, and our ablation study shows clear empirical gains from this repair guidance mechanism.
Importantly, they do not alter the underlying modeling or proof mechanisms and therefore do not affect soundness.
